% Part: computability
% Chapter: computability-theory
% Section: rice-theorem

\documentclass[../../../include/open-logic-section]{subfiles}

\begin{document}

\olfileid{cmp}{thy}{rce}
\olsection{राइसचे प्रमेय}

विचार केल्यास लक्षात येईल की \olref[tot]{prop:total}
च्या सिद्धतेत $\fn{Tot}$ चे विशिष्ट गुणधर्म वापरलेले नाहीत.
$x$ हा $K$ मध्ये असेल तेव्हाच $h(x,y)$ हे स्थिर फलन
$j(y) = 0$ प्रमाणे वागेल अशी रचना आपण केली; पण त्या
परिस्थितीत ते इतर कोणत्याही आंशिक संगणनक्षम फलनाप्रमाणे
वागेल अशीही रचना करता आली असती. या निरीक्षणामुळे अधिक
सर्वसाधारण प्रमेय मांडता येते. साधारणपणे ते सांगते की
संगणनक्षम फलनांचा कोणताही अतुच्छ गुणधर्म निर्णेय नसतो.

$\cfind{0}$, $\cfind{1}$, $\cfind{2}$,~\dots हे आंशिक
संगणनक्षम फलनांचे आपले प्रमाण प्रगणन आहे, हे लक्षात ठेवा.

\begin{thm}[राइसचे प्रमेय]
  $C$ हा आंशिक संगणनक्षम फलनांचा कोणताही संच असू द्या आणि $A =
  \Setabs{n}{\cfind{n} \in C}$ असू द्या. $A$ संगणनक्षम असेल,
  तर $C$ रिकामा असतो किंवा $C$ हा सर्व आंशिक संगणनक्षम
  फलनांचा संच असतो.
\end{thm}

{\em निर्देशांकसंच} म्हणजे असा संच $A$ की $n$ आणि $m$ हे
एकच फलन ``संगणित'' करणारे निर्देशांक असतील, तर $n$ आणि
$m$ हे दोन्ही $A$ मध्ये असतात किंवा दोन्ही नसतात. प्रमेयातील संच~$A$
हा गुणधर्म पूर्ण करतो हे पाहणे कठीण नाही. उलट, $A$
निर्देशांकसंच असेल आणि $C$ हा त्या निर्देशांकांनी संगणित
होणाऱ्या फलनांचा संच असेल, तर $A = \Setabs{n}{\cfind{n} \in C}$.

\begin{explain}
या संज्ञेत राइसचे प्रमेय म्हणजे कोणताही अतुच्छ निर्देशांकसंच
निर्णेय नसतो असे विधान. प्रमेयाचा अर्थ समजण्यासाठी
\emph{कार्यक्रम} (उदा., तुमच्या आवडत्या प्रोग्रामिंग भाषेतील)
आणि ते संगणित करत असलेली फलने यांतील फरक लक्षात घ्या.
कार्यक्रम हे विन्यासात्मक वस्तू आहेत, आणि त्यांच्याविषयीचे
काही प्रश्न नक्कीच संगणनक्षम आहेत: या कार्यक्रमात 150
पेक्षा अधिक चिन्हे आहेत का? त्यात 22 पेक्षा अधिक ओळी
आहेत का? त्यात ``while'' विधान आहे का? ``print'' विधानाचे
आदान म्हणून ``hello world'' ही चिन्हमाला कुठे येते का?
राइसचे प्रमेय सांगते की कार्यक्रमाच्या \emph{वर्तनाविषयीचा}
कोणताही अतुच्छ प्रश्न संगणनक्षम नसतो. उदाहरणार्थ: हा
कार्यक्रम आदान $0$ वर थांबतो का? तो कधी तरी थांबतो का?
तो कधी सम संख्या परत करतो का?
\end{explain}

\begin{proof}[राइसच्या प्रमेयाची सिद्धता]
$C$ रिकामा नाही आणि सर्व आंशिक संगणनक्षम फलनांचा संचही
नाही असे समजा. $A$ हा~$C$ मधील फलनांच्या निर्देशांकांचा
संच असू द्या. $A$ संगणनक्षम असता, तर थांबण्याची समस्या
सोडवता आली असती, हे आपण दाखवू; म्हणून $A$ संगणनक्षम नाही.

व्यापकता न गमावता, कुठेही परिभाषित नसलेले फलन $f$ हे
$C$ मध्ये नाही असे धरू शकतो (तसे नसेल, तर पुढील
युक्तिवादात $C$ आणि त्याचा पूरक संच यांची अदलाबदल करा).
$g$ हे~$C$ मधील कोणतेही फलन असू द्या. कल्पना अशी आहे:
$A$ निर्णीत करता आले, तर~$f$ संगणित करणारे निर्देशांक
आणि~$g$ संगणित करणारे निर्देशांक यांतील फरक ओळखता
येईल; आणि ती क्षमता वापरून थांबण्याची समस्या सोडवता येईल.

कसे ते पाहू. सार्वत्रिक आंशिक संगणनक्षम फलने वापरून
पुढील फलन व्याख्यायित करता येते:
\[
h(x,y) \simeq
\begin{cases}
\text{अपरिभाषित} & \text{जर $\cfind{x}(x) \fundefined$ असेल तर} \\
g(y) & \text{अन्यथा.}
\end{cases}
\]
$h$ संगणित करण्यासाठी आधी $\cfind{x}(x)$ संगणित करण्याचा
प्रयत्न करतो; ते संगणन थांबले, तर पुढे~$g(y)$ संगणित करतो;
आणि {\em हे} संगणन थांबले, तर त्याचे मूल्य परत करतो.
अधिक औपचारिकपणे,
\[
h(x,y) \simeq \Proj{2}{0}(g(y),\fn{Un}(x,x)).
\]
असे लिहिता येते. येथे $\Proj{2}{0}(z_0, z_1) = z_0$ हे
$0$ वे आदान परत करणारे, संगणनक्षम $2$-स्थानी प्रक्षेपण
फलन आहे.

मग $h$ हे आंशिक संगणनक्षम फलनांचे संयोजन आहे. उजवी
बाजू~$g(y)$ इतकीच परिभाषित होते जर आणि फक्त जर
$\fn{Un}(x,x)$ आणि $g(y)$ ही दोन्ही परिभाषित असतील.

निश्चित~$x$ साठी, $\cfind{x}(x)$ अपरिभाषित असेल, तर
प्रत्येक~$y$ साठी $h(x,y)$ अपरिभाषित असते; आणि
$\cfind{x}(x)$ परिभाषित असेल, तर $h(x,y) \simeq g(y)$.
म्हणून निश्चित~$x$ साठी $h(x,y)$ हे एकतर $f$ प्रमाणे
किंवा $g$ प्रमाणे वागते. $\cfind{x}(x)$ परिभाषित आहे
की नाही हे ठरवणे म्हणजे या दोन स्थितींपैकी कोणती आहे
हे ठरवणे. पण हेच $h_x(y) \simeq h(x,y)$ हे~$C$
मध्ये आहे की नाही हे ठरवण्यासारखे आहे; $A$ संगणनक्षम
असता, तर तसे करता आले असते.

अधिक औपचारिकपणे, $h$ आंशिक संगणनक्षम असल्याने तो काही
निर्देशांक~$e$ साठी $\cfind{e}$ या फलनाइतकाच आहे.
$s$-$m$-$n$ प्रमेयानुसार असे आदिम पुनरावर्ती फलन~$s$
आहे की प्रत्येक~$x$ साठी $\cfind{s(e,x)}(y) = h_x(y)$.
आता प्रत्येक $x$ साठी $\cfind{x}(x) \fdefined$ असेल,
तर $\cfind{s(e,x)}$ हे~$g$ इतकेच फलन आहे; म्हणून
$s(e,x)$ हा $A$ मध्ये आहे. उलट, $\cfind{x}(x)
\uparrow$ असेल, तर $\cfind{s(e,x)}$ हे~$f$ इतकेच फलन
आहे; म्हणून $s(e,x)$ हा~$A$ मध्ये नाही. दुसऱ्या शब्दांत,
प्रत्येक~$x$ साठी $x
\in K$ असेल तेव्हाच $s(e,x) \in A$. $A$ संगणनक्षम
असता, तर~$K$ देखील संगणनक्षम ठरला असता; हा विरोधाभास
आहे. म्हणून $A$ संगणनक्षम नाही.
\end{proof}

राइसचे प्रमेय अतिशय प्रभावी आहे. पुढील तत्काळ अनुप्रमेयात
त्याच्या उपयोगांची काही उदाहरणे आहेत.
\begin{cor}
पुढील संच अनिर्णेय आहेत.
\begin{enumerate}
\item $\Setabs{x}{\text{$17$ हे $\cfind{x}$ च्या व्याप्तीत आहे}}$
\item $\Setabs{x}{\text{$\cfind{x}$ स्थिर आहे}}$
\item $\Setabs{x}{\text{$\cfind{x}$ सर्वत्र परिभाषित आहे}}$
\item $\Setabs{x}{\text{जेव्हा $y < y'$ असेल तेव्हा $\cfind{x}(y) \fdefined$, आणि
    जर $\cfind{x}(y') \fdefined$ असेल, तर $\cfind{x}(y) < \cfind{x}(y')$}}$
\end{enumerate}
\end{cor}

\begin{proof}
  हे सर्व अतुच्छ निर्देशांकसंच आहेत.
\end{proof}

\end{document}
